%% Capítulo 12: Conclusiones y Trabajo Futuro

\chapter{Conclusiones y Trabajo Futuro}
\label{chap:conclusiones}

\section{Síntesis de Resultados}
\label{sec:sintesis}

\section{Contribuciones de la Tesis}
\label{sec:contribuciones-finales}

\section{Limitaciones}
\label{sec:limitaciones}

Las principales limitaciones de este trabajo son:

\begin{enumerate}
  \item \textbf{Riesgo de contaminación de datos (\textit{data contamination}) en el
  período de evaluación}: el protocolo de benchmark compara los pronósticos del DSGE
  y del sistema de agentes LLM contra datos realizados en el tramo
  2023T1--2024T4 (\texttt{EVAL\_START}/\texttt{EVAL\_END} en
  \texttt{dsge/evaluation.py}). Los modelos utilizados para los agentes
  (\texttt{gpt-4o-mini}, \texttt{claude-haiku-4-5}, \texttt{claude-sonnet-4-6}) tienen
  fechas de corte de entrenamiento posteriores a ese período, y la crisis
  macroeconómica argentina de 2023--2024 (inflación superior al 200\% anual,
  elección presidencial de noviembre de 2023, devaluación post-electoral,
  unificación cambiaria, debate de dolarización) tuvo cobertura mediática
  internacional extensa. Esto implica que los agentes LLM pueden tener acceso, vía
  memoria paramétrica, a los resultados macroeconómicos reales del propio período
  que se les pide pronosticar. El prompt de cada agente incluye además el período
  calendario exacto (\texttt{"Período: \{state.period\}"} en
  \texttt{agents/base.py}), lo cual facilita que el modelo asocie el estado
  simulado con el episodio histórico real en lugar de generar una decisión a partir
  únicamente de las variables de estado provistas. A diferencia del DSGE --cuyo
  pronóstico es una extrapolación genuinamente fuera de muestra a partir de
  parámetros estimados en 2016--2022, sin acceso a los datos de 2023--2024--, no es
  posible garantizar que las decisiones de los agentes LLM constituyan un ejercicio
  de pronóstico genuino y no, parcial o totalmente, una recuperación de información
  memorizada durante el entrenamiento. Esta asimetría amenaza la validez interna de
  la comparación DSGE~vs.~agentes para el tramo evaluado y debería tratarse como una
  limitación central al interpretar cualquier resultado donde los agentes LLM
  igualen o superen al DSGE en las métricas de error.

\end{enumerate}

\section{Trabajo Futuro}
\label{sec:trabajo-futuro}
